% !TeX root = ../main.tex
\begin{frame}{03 创新思路——支架应用构想}
  \begin{SlideCanvas}
    \BodyAt{55}{84}{850}{以课题组反手性结构为候选机制，在给定长度偏差下匹配目标径向力曲线，研究几何与承载要求的共同实现。}
    \PicAt{55}{160}{417}{235}{assets/group_candidate_array.pdf}
    \CaptionAt{55}{404}{417}{课题组已发表单胞与阵列，作为候选机制来源}
    \DetailAt{55}{441}{417}{卷绕成管后，按周向工况重新设计杆件与连接。}
    \PointAt{530}{155}{370}{全过程长度约束}
    \BodyAt{530}{202}{370}{$\displaystyle
      \max_{D\in[D_a,D_b]}\left|\frac{L(D)-L_0}{L_0}\right|\leq\delta_L$}
    \PointAt{530}{295}{370}{径向承载目标}
    \BodyAt{530}{339}{370}{$\displaystyle F_r(D)\approx F_r^\ast(D)$}
    \DetailAt{530}{407}{370}{统一材料量、最小杆宽和应变约束，在同一加载分支比较目标曲线。}
    \Sources{图为\RefGroup，图4(a)(b)，代表已发表几何机制。本页为支架应用设想，未展示新增优化结果。}
  \end{SlideCanvas}
  \BICConclusionBand{应用构想聚焦于\BICBandEmph{长度轨迹与径向力曲线的联合定制}}
  \note{保留个人思考但减少比重，回应课程对自身思考的要求。图件取自课题组已发表结果，不是新支架构型，不用已发表几何代替未完成的优化结果。已有零泊松比支架兼顾短缩与径向性能，本设想研究在特定材料、直径区间与约束下对完整目标曲线的共同反设计。目标径向曲线需由应用及参考器件确定。可以先用既有超弹性程序筛选候选，再将圆管的周向加载响应反馈优化；实际NiTi或球囊金属部署需相应本构。}
\end{frame}
